\subsection{乘积的导数}

现在考虑 \(p\) 个拓扑向量空间 \(\mathrm{E}_i\) （ \(1 \leqslant i \leqslant p\) ，在 \(\mathbf{R}\) 上），以及一个连续多线性 \(^1\) 映射（我们将把它记作

\[
\left(\mathbf {x} _ {1}, \mathbf {x} _ {2}, \dots , \mathbf {x} _ {p}\right) \mapsto \left[ \mathbf {x} _ {1}. \mathbf {x} _ {2} \dots \mathbf {x} _ {p} \right])
\]

它把 \(E_{1} \times E_{2} \times \cdots \times E_{p}\) 映入 R 上的一个拓扑向量空间 F 中。

命题 3. 对每个指标 \(i\) ( \(1 \leqslant i \leqslant p\) )，设 \(\mathbf{f}_i\) 是定义在区间 \(\mathrm{I} \subset \mathbf{R}\) 上、取值于 \(\mathrm{E}_i\) 中、且在点 \(x_0 \in \mathrm{I}\) 可导的函数。则函数

\[
x \mapsto [ \mathbf {f} _ {1} (x). \mathbf {f} _ {2} (x) \dots \mathbf {f} _ {p} (x) ]
\]

定义在 I 上、取值于 F 中，它具有导数，等于

\[
\sum_ {i = 1} ^ {p} \left[ \mathbf {f} _ {1} (x _ {0}) \dots \mathbf {f} _ {i - 1} (x _ {0}). \mathbf {f} _ {i} ^ {\prime} (x _ {0}). \mathbf {f} _ {i + 1} (x _ {0}) \dots \mathbf {f} _ {p} (x _ {0}) \right]\tag{1}
\]

（在点 \(x_0\) 处）。

令 \(\mathbf{h}(x) = [\mathbf{f}_1(x).\mathbf{f}_2(x)\dots \mathbf{f}_p(x)]\) ；则由恒等式

\[
\left[ \mathbf {b} _ {1}. \mathbf {b} _ {2} \dots \mathbf {b} _ {p} \right] - \left[ \mathbf {a} _ {1}. \mathbf {a} _ {2} \dots \mathbf {a} _ {p} \right] = \sum_ {i = 1} ^ {p} \left[ \mathbf {b} _ {1} \dots \mathbf {b} _ {i - 1}. (\mathbf {b} _ {i} - \mathbf {a} _ {i}). \mathbf {a} _ {i + 1} \dots \mathbf {a} _ {p} \right],
\]

我们便可写出

\[
\mathbf {h} (x) - \mathbf {h} (x _ {0}) = \sum_ {i = 1} ^ {p} \left[ \mathbf {f} _ {1} (x) \dots \mathbf {f} _ {i - 1} (x). (\mathbf {f} _ {i} (x) - \mathbf {f} _ {i} (x _ {0})). \mathbf {f} _ {i + 1} (x _ {0}) \dots \mathbf {f} _ {p} (x _ {0}) \right].
\]

把等式两边同乘 \(\frac{1}{x - x_0}\) 并让 \(x\) 在 I 中趋于 \(x_0\) ，便得到表达式 (1)，因为映射

\[
(\mathbf {x} _ {1}, \mathbf {x} _ {2}, \dots , \mathbf {x} _ {p}) \mapsto [ \mathbf {x} _ {1}. \mathbf {x} _ {2} \dots \mathbf {x} _ {p} ]
\]

与 F 中的加法都是连续的。

\(^{1}\) 回忆（Alg., II, p.~265）下面的定义：如果从 \(E_{1} \times E_{2} \times \cdots \times E_{p}\) 到 F 中的映射 f 的每个偏映射

\[
\mathbf {x} _ {i} \mapsto \mathbf {f} (\mathbf {a} _ {1}, \dots , \mathbf {a} _ {i - 1}, \mathbf {x} _ {i}, \mathbf {a} _ {i + 1}, \dots , \mathbf {a} _ {p})
\]

从 \(E_{i}\) 到 F （ \(1 \leqslant i \leqslant p\) ）都是线性映射，其中 \(a_{j}\) （对于 \(j \neq i\) 的下标）在 \(E_{j}\) 中任意取定，则称 f 是多线性的。我们指出，若诸 \(E_{i}\) 在 R 上有限维，则从 \(E_{1} \times E_{2} \times \cdots \times E_{p}\) 到 F 中的每个多线性映射必定连续；当这些空间中有一些是无穷维拓扑向量空间时，则未必如此。

当某些函数 \(\mathbf{f}_i\) 为常数时，表达式 (1) 中含有其导数 \(\mathbf{f}_i'(x_0)\) 的项为零。

我们详细考察应用中最重要的特例 \(p = 2\) ：若 \((\mathbf{x}, \mathbf{y}) \mapsto [\mathbf{x}.\mathbf{y}]\) 是 \(\mathrm{E} \times \mathrm{F}\) 到 \(\mathrm{G}\) 中的连续双线性映射（E、F、G 为 \(\mathbf{R}\) 上的拓扑向量空间），而 \(\mathbf{f}\) 和 \(\mathbf{g}\) 是两个向量函数，分别在点 \(x_0\) 可导，各取值于 E 和 F 中，则向量函数 \(x \mapsto [\mathbf{f}(x).\mathbf{g}(x)]\) （我们把它记作 \([\mathbf{f}.\mathbf{g}]\) ）具有导数，等于 \([\mathbf{f}'(x_0).\mathbf{g}(x_0)] + [\mathbf{f}(x_0).\mathbf{g}'(x_0)]\) （在 \(x_0\) 处）。特别地，若 \(\mathbf{a}\) 是常向量，则 \([\mathbf{a}.\mathbf{f}]\) （相应地，\([\mathbf{f}.\mathbf{a}]\) ）具有导数，等于 \([\mathbf{a}.\mathbf{f}'(x_0)]\) （相应地，\([\mathbf{f}'(x_0).\mathbf{a}]\) ）（在 \(x_0\) 处）。

若 \(\mathbf{f}\) 和 \(\mathbf{g}\) 都在 I 上可导，则 {[}f.g{]} 也在 I 上可导，且我们有

\[
\left[ \mathbf {f}. \mathbf {g} \right] ^ {\prime} = \left[ \mathbf {f} ^ {\prime}. \mathbf {g} \right] + \left[ \mathbf {f}. \mathbf {g} ^ {\prime} \right].\tag{2}
\]

例。1) 设 \(f\) 是实函数，\(\mathbf{g}\) 是向量函数，二者都在点 \(x_0\) 可导；函数 \(\mathbf{g}f\) 具有导数，等于 \(\mathbf{g}'(x_0)f(x_0) + \mathbf{g}(x_0)f'(x_0)\) （在 \(x_0\) 处）。特别地，若 \(\mathbf{a}\) 为常数，则 \(\mathbf{a}f\) 具有导数 \(\mathbf{a}f'(x_0)\) 。最后这个注记结合 I, p.~5 的例 1 证明：若 \(\mathbf{f} = (f_i)_{1 \leqslant i \leqslant n}\) 是取值于 \(\mathbf{R}^n\) 中的向量函数，则 \(\mathbf{f}\) 在点 \(x_0\) 可导当且仅当每个实函数 \(f_i\) ( \(1 \leqslant i \leqslant n\) ) 都在该点可导：因为若 \((\mathbf{e}_i)_{1 \leqslant i \leqslant n}\) 是 \(\mathbf{R}^n\) 的规范基，我们便可写 \(\mathbf{f} = \sum_{i=1}^{n} \mathbf{e}_i f_i\) 。

\begin{enumerate}
\def\labelenumi{\arabic{enumi})}
\setcounter{enumi}{1}
\tightlist
\item
  实函数 \(x^{n}\) 得自多线性函数
\end{enumerate}

\[
(x _ {1}, x _ {2}, \ldots , x _ {n}) \mapsto x _ {1} x _ {2} \ldots x _ {n}
\]

它定义在 \(R^{n}\) 上，是通过把每个 \(x_{i}\) 都换成 x 而得到的；于是命题 3 表明 \(x^{n}\) 在 R 上可导且具有导数 \(nx^{n-1}\) 。由此，多项式函数 \(a_{0}x^{n} + a_{1}x^{n-1} + \cdots + a_{n-1}x + a_{n}\) （诸 \(a_{i}\) 为常向量）具有导数

\[
n \mathbf {a} _ {0} x ^ {n - 1} + (n - 1) \mathbf {a} _ {1} x ^ {n - 2} + \dots + \mathbf {a} _ {n - 1};
\]

当诸 \(\mathbf{a}_i\) 为实数时，此函数与代数学（A, IV）中所定义的多项式函数的导数相一致。

\begin{enumerate}
\def\labelenumi{\arabic{enumi})}
\setcounter{enumi}{2}
\tightlist
\item
  欧氏数量积 \((\mathbf{x}|\mathbf{y})\) (Gen.~Top., VI, p.~40) 是 \(R^{n}\times R^{n}\) 到 R 中的一个双线性映射（必然连续）。若 f 和 g 是取值于 \(R^{n}\) 、且在点 \(x_{0}\) 可导的两个向量函数，则实函数 \(x\mapsto(\mathbf{f}(x)|\mathbf{g}(x))\) 具有导数，等于 \((\mathbf{f}'(x_{0})|\mathbf{g}(x_{0}))+(\mathbf{f}(x_{0})|\mathbf{g}'(x_{0}))\) （在点 \(x_{0}\) 处）。对 \(C^{n}\) 上的埃尔米特数量积有类似的结果，这时把该空间看作 R 上的向量空间。
\end{enumerate}

我们特别考察欧氏范数 \(\|\mathbf{f}(x)\|\) 为常数的情形，此时 \((\mathbf{f}(x) \mid \mathbf{f}(x)) = \|\mathbf{f}(x)\|^{2}\) 也是常数；写出 \((\mathbf{f}(x) \mid \mathbf{f}(x))\) 的导数在 \(x_{0}\) 处为零，便得到 \((\mathbf{f}(x_{0}) \mid \mathbf{f}'(x_{0})) = 0\) ；换句话说，\(\mathbf{f}'(x_{0})\) 与 \(\mathbf{f}(x_{0})\) 正交。

\begin{enumerate}
\def\labelenumi{\arabic{enumi})}
\setcounter{enumi}{3}
\item
  若 E 是 R 上的一个拓扑代数（cf.~Introduction），E 的两个元素之积 xy 便是 \((\mathbf{x}, \mathbf{y})\) 的连续双线性函数；若 f 与 g 取值于 E 且在点 \(x_{0}\) 可导，则函数 \(x \mapsto \mathbf{f}(x)\mathbf{g}(x)\) 具有导数，等于 \(\mathbf{f}'(x_{0})\mathbf{g}(x_{0}) + \mathbf{f}(x_{0})\mathbf{g}'(x_{0})\) （在 \(x_{0}\) 处）。特别地，若 \(U(x) = (\alpha_{ij}(x))\) 与 \(V(x) = (\beta_{ij}(x))\) 是两个 n 阶方阵，在 \(x_{0}\) 可导，则其乘积 UV 具有导数，等于 \(U'(x_{0})V(x_{0}) + U(x_{0})V'(x_{0})\) （在 \(x_{0}\) 处）（其中 \(U'(x) = (\alpha'_{ij}(x))\) ，\(V'(x) = (\beta'_{ij}(x)))\) 。
\item
  行列式 \(\operatorname{det}(\mathbf{x}_1, \mathbf{x}_2, \ldots, \mathbf{x}_n)\) （\(n\) 个向量 \(\mathbf{x}_i = (x_{ij})_{1 \leqslant j \leqslant n}\) 的行列式，这些向量取自空间 \(\mathbf{R}^n\) ；Alg., III, p.~522）既然是诸 \(\mathbf{x}_i\) 的（连续）多线性函数，便可以看出：若 \(n^2\) 个实函数 \(f_{ij}\) 都在 \(x_0\) 可导，则它们的行列式 \(g(x) = \det(f_{ij}(x))\) 具有导数，等于
\end{enumerate}

\[
\sum_ {i = 1} ^ {n} \left[ \mathbf {f} _ {1} (x _ {0}), \dots , \mathbf {f} _ {i - 1} (x _ {0}), \mathbf {f} _ {i} ^ {\prime} (x _ {0}), \mathbf {f} _ {i + 1} (x _ {0}), \dots , \mathbf {f} _ {n} (x _ {0}) \right]
\]

（在 \(x_{0}\) 处），其中 \(\mathbf{f}_{i}(x)=(f_{ij}(x))_{1\leqslant j\leqslant n}\) ；换句话说，对每个 i，把 n 阶行列式中第 \(i^{th}\) 列的元素换成它们的导数，这样形成的 n 个行列式之和便是该行列式的导数。

注. 若 \(U(x)\) 是在点 \(x_{0}\) 可导且可逆的方阵，则其行列式 \(\Delta(x) = \det(U(x))\) 的导数可通过 \(U(x)\) 的导数由公式

\[
\varDelta^ {\prime} (x _ {0}) = \varDelta (x _ {0}). \mathrm{Tr} (U ^ {\prime} (x _ {0}) U ^ {- 1} (x _ {0})).\tag{3}
\]

表出。事实上，令 \(U(x_0 + h) = U(x_0) + hV\) ；则按定义，\(V\) 趋于 \(U'(x_0)\) （当 \(h\) 趋于 0 时）。可以写出

\[
\Delta (x _ {0} + h) = \Delta (x _ {0}). \det (I + h V U ^ {- 1} (x _ {0})).
\]

现在 \(\det(I + hX) = 1 + h\operatorname{Tr}(X) + \sum_{k=2}^{n} \lambda_k h^k\) ，其中 \(\lambda_k (k \geqslant 2)\) 是矩阵 X 的元素的多项式；由于当 h 趋于 0 时 \(VU^{-1}(x_0)\) 的元素有极限，我们确实得到公式 (3)。

